\section{Discussion}
\label{sec:analysis}

\subsection{Why Scan Recovery Changes the Comparison}
Accepted LiDAR scans repeatedly constrain pose, so gravity-state drift need
not produce a worse trajectory. During an outage, errors accumulate until
registration resumes from the propagated state.

For a small gravity-direction discrepancy $\varepsilon$ propagated without an
external update for time $T$, the leading position contribution is
\begin{equation}
 \|\delta\mathbf p_g\|\approx\tfrac12g\sin(\varepsilon)T^2 .
 \label{eq:timescale}
\end{equation}
This describes displacement before recovery, not final trajectory error.
Registration may correct a large error but leave a smaller one partly
uncorrected, depending on the scene and returning pose. This is consistent
with the non-monotonic 3-s and 5-s results, although we do not directly measure
the registration convergence basin.

In the matched switch, both runs start from the same state and map, but fixing
gravity conditions the covariance. Gravity has no process dynamics, so the
nominal means remain equal during the first gap; Eq.~\ref{eq:timescale} cannot
explain separation before scan return. Differences arise during correction,
through the changed covariance and continued Online gravity updates. The APE
traces show accumulation over successive recovery corrections. ATE
interactions stay positive across tested starts, while RMSE$_z$ interactions
change sign. The repeated advantage concerns 3D recovery, not direct height
information.

The relevant interval $T_k=t_k-t_{k-1}$ is between \emph{accepted pose
corrections}, not arriving LiDAR packets. A rejected scan cannot interrupt
inertial error accumulation. Message removal controls these intervals, whereas
natural rejection may depend on geometry and the estimated pose. Biased but
accepted registrations instead introduce erroneous corrections, a different
failure mode not tested here. Scan rate alone is therefore insufficient to
transfer the observed thresholds.

\subsection{Interpreting the Direction-Factor Response}
The factor uses AHRS directions derived from the IMU measurements also used
for preintegration. It observes neither height nor vertical velocity, so its
vertical-error effects are mediated by attitude, bias, propagation, and
registration.

At fixed gravity magnitude, a direction error $\varepsilon$ gives
\begin{equation}
 \begin{aligned}
  \|\mathbf e_g^\top\delta\mathbf g\|
    &=g(1-\cos\varepsilon)=\mathcal O(\varepsilon^2),\\
  \|\mathbf P_{\perp}\delta\mathbf g\|
    &=g\sin\varepsilon=\mathcal O(\varepsilon).
 \end{aligned}
\label{eq:gravity_components}
\end{equation}
where $\mathbf e_g$ is unit true down and
$\mathbf P_\perp=\mathbf I-\mathbf e_g\mathbf e_g^\top$.
Small direction errors therefore act to first order transversely, but only to
second order along true down. A direction factor can still change RMSE$_z$ by
rotating measured specific force or changing scan-matching initialization.
Neither route measures height, and Eq.~\ref{eq:gravity_components} alone does
not determine the final error change.

Hall05's improvements and TUHH's mixed or adverse responses reject an
unconditional benefit from stronger direction consistency. They do not
invalidate independent gravity, velocity, contact, radar, or registration
information
\cite{kubelka2022gravity,noh2025garlio,noh2025garlileo}, nor do they test a
downstream global pose-graph prior. Factor covariance must be validated with
the state structure, correction regime, and both trajectory metrics.

Unmodeled same-IMU correlation could contribute to the adverse response, but
we do not isolate it from weighting and online-gravity interactions.
A smaller residual establishes agreement with AHRS, not an independent
reference (Fig.~\ref{fig:direction}). Table~\ref{tab:direction} also retains
resets and arc warnings; reporting only successful accuracy comparisons would
overstate reliability.

\begin{table}[t]
\caption{Direction-factor effects relative to factor-off.}
\label{tab:direction}
\centering
\footnotesize
\setlength{\tabcolsep}{2.0pt}
% Auto-generated by scripts/make_tables.py; do not edit.
\begingroup
\renewcommand{\arraystretch}{1.05}
\begin{tabular*}{\columnwidth}{@{\extracolsep{\fill}}llrrl@{}}
\toprule
Weight & $g$ state & $\Delta$RMSE$_z$ [\%] & $\Delta$ATE [\%] & outcome \\
\midrule
\rowcolor{TableBand}
\multicolumn{5}{@{}l}{\textsc{Hall05}\enspace 5-s dropout per 20 s} \\
$0.5^\circ$ & \methodG{} & $-42.4$ & $-76.5$ & 3/3 \\
$2^\circ$ & \methodG{} & $-43.8$ & $-75.0$ & 3/3 \\
$5^\circ$ & \methodG{} & $-30.5$ & $-69.1$ & 2/3; 1 R \\
$0.5^\circ$ & \methodA{} & $-9.6$ & $-11.2$ & 3/3 \\
$2^\circ$ & \methodA{} & $-7.6$ & $-5.2$ & 3/3 \\
$5^\circ$ & \methodA{} & $-2.3$ & $-4.0$ & 3/3 \\
\addlinespace[2pt]
\rowcolor{TableBand}
\multicolumn{5}{@{}l}{\textsc{TUHH handheld}\enspace 5-s dropout per 20 s} \\
$0.5^\circ$ & \methodG{} & $-64.6$ & $+97.6$ & 3/3; 1 A \\
$2^\circ$ & \methodG{} & $+56.2$ & $+25.0$ & 3/3 \\
$5^\circ$ & \methodG{} & $-38.3$ & $-48.1$ & 3/3 \\
$0.5^\circ$ & \methodA{} & $+21.8$ & $+230.9$ & 3/3; 1 A \\
$2^\circ$ & \methodA{} & $+450.0$ & $+108.4$ & 3/3 \\
$5^\circ$ & \methodA{} & $+65.4$ & $+8.2$ & 3/3 \\
\bottomrule
\end{tabular*}
\endgroup

\vspace{1pt}
{\scriptsize\raggedright Outcome entries report accuracy runs/planned runs;
R denotes a graph reset and A a retained arc warning.\par}
\end{table}

\subsection{Choosing a Default Configuration}
Keeping $\mathbf g$ and $\mathbf b_a$ online allows adjustment but does not
guarantee physical accuracy. The handheld runs retain different gravity and
bias estimates despite similar effective accelerations and poses. Vehicle
excitation largely removes those differences. As Eq.~\ref{eq:coupling}
suggests, accurate pose does not establish that gravity and bias have each
been identified.

With reliable initialization and continuous correction, FAST-LIO2 FixG is
equivalent on average but saves negligible compute. Keeping both states online
allows adaptation when motion contaminates startup or corrections disappear.
Fixed-bias variants are phase-sensitive, and one FixG start incurs large ATE
harm. We therefore recommend keeping both states online by default. FixG is an
option for a validated steady-state setting, not a generally better estimator.
A separate direction factor needs its own paired RMSE$_z$ and ATE validation
under the intended operating conditions.

Across \NumRuntimePairs{} serial pairs, FixG reduced filter algebra by
\RuntimeAlgebraReduction{}, but it occupied only \RuntimeAlgebraCoreShare{} of
core time versus \RuntimeMatchingCoreShare{} for scan matching. Its attributable
saving was \RuntimeAttributableCoreSaving{}, and whole-core change
(\RuntimeCoreMedianChangeRange{}) crossed zero. Dimension reduction is not a
runtime strategy.

Two dropout and two initialization trajectories cannot establish universal
thresholds. The factor study retains \NumHallDirResetOutcomes{} reset and
\NumTuhhDirArcWarnings{} arc warnings. Larger initialization errors,
calibration faults, and motion-aware initialization also remain untested.
TOST concerns pose error, not covariance
consistency; dropout results are not pooled across trajectories.

% balance.sty and ieeeconf can differ by a few points in the final output box;
% this tolerance suppresses that rounding-only diagnostic without changing layout.
\vfuzz=3pt
\balance
